\begin{abstract}
Η παρούσα διπλωματική εργασία πραγματεύεται τον σχεδιασμό, την υλοποίηση και την αξιολόγηση μιας ολοκληρωμένης πλατφόρμας Ψηφιακού Διδύμου (\en{Digital Twin}), ειδικά προσαρμοσμένης για εκπαιδευτικά περιβάλλοντα και Έξυπνες Πανεπιστημιουπόλεις (\en{Smart Campuses}). Σε αντίθεση με τις παραδοσιακές, αντιδραστικές (\en{reactive}) μεθόδους διαχείρισης κτιρίων ή τα απλά συστήματα παρακολούθησης (\en{Digital Shadows}), η προτεινόμενη λύση εγκαθιδρύει μια αμφίδρομη, πραγματικού χρόνου επικοινωνία μεταξύ του φυσικού σχολικού χώρου και του ψηφιακού του αντιγράφου, επιτρέποντας τη λήψη αυτοματοποιημένων αποφάσεων.

Η ανάπτυξη του συστήματος ακολούθησε μια αυστηρή διεπιστημονική προσέγγιση, βασισμένη στην αρχιτεκτονική πέντε επιπέδων. Σε επίπεδο υλικού, εγκαταστάθηκε δίκτυο κόμβων Διαδικτύου των Πραγμάτων (\en{IoT}) χαμηλού κόστους, αποτελούμενο από μικροελεγκτές ESP32, έξυπνους μετρητές Shelly 1PM και συσκευές Sonoff. Σε επίπεδο λογισμικού, αναπτύχθηκε μια ισχυρή, τοπικά εκτελούμενη αρχιτεκτονική με χρήση του πρωτοκόλλου MQTT, βάσης χρονοσειρών InfluxDB και του ανοιχτού κώδικα ενδιάμεσου λογισμικού openHAB, η οποία τροφοδοτεί μια προσαρμοσμένη τρισδιάστατη (\en{3D}) διεπαφή χρήστη τεχνολογίας React.

Η έρευνα δομείται πάνω σε τρεις κεντρικούς πυλώνες: (α) τη βελτιστοποίηση του Μαθησιακού Περιβάλλοντος (\en{Environmental Quality}), μέσω της διαρκούς παρακολούθησης της ακουστικής άνεσης για την ενίσχυση της συγκέντρωσης των μαθητών βάσει των προτύπων του ΠΟΥ (όριο 55 dBA), (β) την Ενεργειακή Αποδοτικότητα (\en{Energy Efficiency}), με συνεχή καταγραφή ισχύος, και (γ) τη Λειτουργική Διαχείριση (\en{Facility Management}) μέσω απομακρυσμένου ελέγχου. Από την πιλοτική εφαρμογή στα «Εκπαιδευτήρια Πλάτων» (Μάρτιος – Νοέμβριος 2025), επιτεύχθηκε και επαληθεύτηκε μέση μείωση της ενεργειακής κατανάλωσης κατά 19,7%. Παράλληλα, η τεχνοοικονομική ανάλυση επιβεβαίωσε την υψηλή προστιθέμενη αξία της μετάβασης σε ένα σύγχρονο Έξυπνο Σχολείο, καταδεικνύοντας απόσβεση της επένδυσης σε 42,5 μήνες και 5ετές ROI 41,02%.

\textbf{Λέξεις-Κλειδιά:} Ψηφιακό Δίδυμο, Έξυπνο Σχολείο, Διαδίκτυο των Πραγμάτων (\en{IoT}), Μαθησιακό Περιβάλλον, Ενεργειακή Αποδοτικότητα, Διαχείριση Εγκαταστάσεων.
\end{abstract}

\renewcommand{\abstractname}{Abstract}
\begin{abstract}
This diploma thesis presents the design, implementation, and evaluation of a comprehensive \en{Digital Twin} platform, tailored specifically for educational environments and \en{Smart Campuses}. In contrast to traditional, reactive building management methods or simple monitoring systems (\en{Digital Shadows}), the proposed solution establishes a bidirectional, real-time communication between the physical school space and its digital replica, enabling automated decision-making.

The system's development followed a rigorous interdisciplinary approach, based on a five-layer architecture. At the hardware level, a low-cost \en{Internet of Things} (\en{IoT}) node network was deployed, consisting of ESP32 microcontrollers, Shelly 1PM smart meters, and Sonoff devices. At the software level, a robust, locally executed architecture was developed using the MQTT protocol, InfluxDB time-series database, and the open-source openHAB middleware, which feeds a custom three-dimensional (\en{3D}) React-based user interface.

The research is structured around three core pillars: (\en{a}) the optimization of the Learning Environment (\en{Environmental Quality}), through continuous acoustic comfort monitoring to enhance student concentration based on WHO standards (\en{55 dBA threshold}), (\en{b}) Energy Efficiency, with continuous power tracking, and (\en{c}) \en{Facility Management} via remote control. From the pilot implementation at "Platon Schools" (\en{March – November 2025}), an average energy consumption reduction of 19.7% was achieved and verified. Concurrently, the techno-economic analysis confirmed the high added value of transitioning to a modern \en{Smart Campus}, demonstrating an investment payback period of 42.5 months and a 5-year ROI of 41.02%.

\textbf{Keywords:} \en{Digital Twin}, \en{Smart Campus}, \en{Internet of Things} (\en{IoT}), Learning Environment, Energy Efficiency, \en{Facility Management}.
\end{abstract}

